\section{Unitary Representations of Continuous Symmetries}

If the symmetry group G is a Lie group, then the quantum unitary operators representing the group action on states and observables can be expressed as functions of the parameter vector \(\underline{\theta}\) of \(G\):

[...]

\[
    U(m(\theta, \theta')) = U(g(m(\theta, \theta'))) = U(g(\theta) g(\theta')) = \gamma(g(\theta), g(\theta')) U(g(\theta)) U(g(\theta')),
\]
where \(\gamma(g(\theta), g(\theta'))\) is a phase factor which can be written as \(\gamma(g(\theta), g(\theta')) = e^{i \omega(g(\theta), g(\theta'))}\) for some real function \(\omega\). Thus
\[
    U(m(\theta, \theta')) = e^{i \omega(g(\theta), g(\theta'))} U(g(\theta)) U(g(\theta')).
\]

\begin{example}[EXCERCISE]

    ??? photo

\end{example}

When you have a Lie Group and a projected unitary representation in the Hilbert space of states, you can always get to this relation with the phase factor \(\gamma\) by redefining the unitary operators \(U(g)\) with a suitable phase factor, all given to the properties of the Lie group \(G\) and the properties of the projective representation \(U\).

Since the Lie groups act homogeneously, we can always translates what happens at the identity element to any other element of the group. Thus, by looking at the properties of the group around the identity element, we can understand the properties of the whole group. By Taylor expanding \(U(\underline{\theta})\) near the neutral element \(\underline{0}\) of the group, we get
\[
    U(\underline{\theta}) = 1 - i \theta^a Q_a + \frac{1}{2} \theta^a \theta^b Q_{ab} + \mathrm{O}(\theta^3), \quad Q_{ab} = Q_{ba},
\]
where \(Q_a\) are the generators of the representation \(U\) of the Lie group \(G\). The generators \(Q_a\) are essentially the tangent vectors to the representation \(U\) of the Lie group \(G\) at the identity element
\[
    Q_a = i \frac{\partial U(\underline{\theta})}{\partial \theta^a} \Big|_{\underline{\theta} = \underline{0}}.
\]

We are now going to prove that the generators \(Q_a\) of the representation \(U\) of the Lie group \(G\) are self-adjoint operators:
\[
    Q_a^{\dagger} = Q_a.
\]
We can start from the adjoint
\[
    \begin{aligned}
        Q_a^{\dagger} & = \left(i \frac{\partial U(\underline{\theta})}{\partial \theta^a} \Big|_{\underline{\theta} = \underline{0}}\right)^{\dagger}                                        \\
                      & = -i \frac{\partial U^{\dagger}(\underline{\theta})}{\partial \theta^a} \Big|_{\underline{\theta} = \underline{0}}                                                    \\
                      & = -i \frac{\partial U^{-1}(\underline{\theta})}{\partial \theta^a} \Big|_{\underline{\theta} = \underline{0}}                                                         \\
                      & = +i U(\underline{\theta})^{-1} \frac{\partial U(\underline{\theta})}{\partial \theta^a} U(\underline{\theta})^{-1} \Big|_{\underline{\theta} = \underline{0}} = Q_a,
    \end{aligned}
\]
where we have used the unitarity of \(U\) and the fact that \(U(\underline{0}) = 1\). [comments]

We have another property to prove: we have to normalize the phase function \(\omega\)
\[
    \omega(\underline{\theta}, \underline{0}) = \omega(\underline{0}, \underline{\theta}) = 0,
\]
so that we remember that
\[
    \omega(\underline{\theta}, \underline{\theta}') = \theta^a \theta'^b \omega_{ab} + \mathrm{O}(\theta^3),
\]
so that the phase factor \(\gamma\) is normalized to 1 at the identity element of the group, since
\[
    e^{i \omega(\underline{0}, \underline{\theta})} = \gamma(g(\underline{0}), g(\underline{\theta}')) = 1 = e^{i \omega(\underline{\theta}, \underline{0})} = \gamma(g(\underline{\theta}), g(\underline{0})) = 1.
\]
???

[...]

\[
    \left[ Q_a, Q_b \right]
\]

...
as we did for \(g(\theta)g(\theta')\) we can multiply the taylor expansion of \(U(\theta)\) and \(U(\theta')\) and compare the result with the taylor expansion of \(U(m(\theta, \theta'))\) to get
\[
    U(\underline{\theta}) U(\bs{\theta'}) = e^{-i \omega(\underline{\theta}, \bs{\theta'})} U(m(\underline{\theta}, \bs{\theta'})).
\]
If we insert the taylor exapnsions of
\begin{itemize}
    \item \(U(\underline{\theta}) = 1 - i \theta^a Q_a + \frac{1}{2} \theta^a \theta^b Q_{ab} + \mathrm{O}(\theta^3)\)
    \item \(\omega(\underline{\theta}, \bs{\theta'}) = \theta^a \theta'^b \omega_{ab} + \mathrm{O}(\theta^3)\)
    \item \(m(\underline{\theta}, \bs{\theta'}) = \underline{\theta} + \bs{\theta'} + \frac{1}{2} f_{ab}^c \theta^a \theta'^b + \mathrm{O}(\theta^3)\)
\end{itemize}
we get a mess, but if we compute everything and compare the coefficients of the terms of order \(\theta^a \theta'^b\) we get
\[
    Q_a Q_b = i m^c_{ab} Q_c + i \omega_{ab}\mathbb{I} - Q_{ab},
\]
and since we wanted to compute the commutator \(\left[ Q_a, Q_b \right]\) we can compute
\[
    Q_b Q_a = i m^c_{ba} Q_c + i \omega_{ba}\mathbb{I} - Q_{ba},
\]
so that subtracting the two we get
\[
    \left[ Q_a, Q_b \right] = i (m^c_{ab} - m^c_{ba}) Q_c + i (\omega_{ab} - \omega_{ba})\mathbb{I} = i f^c_{ab} Q_c + i \kappa_{ab} \mathbb{I},
\]
where we have recalled the structure constants of the Lie algebra \(\mathfrak{g}\) of the Lie group \(G\) and we have defined \(\kappa_{ab} = \omega_{ab} - \omega_{ba}\). This last coefficients \(\kappa_{ab}\) are called the \textbf{central charges} of the representation \(U\) of the Lie group \(G\). A central algebra is an algebra where the central charges are zero, so that the commutator of the generators \(Q_a\) of the representation \(U\) of the Lie group \(G\) is just
\[
    \left[ Q_a, Q_b \right] = i f^c_{ab} Q_c.
\]
The properties of the commutators of the generators \(Q_a\) of the representation \(U\) can be expressed as
\begin{itemize}
    \item \(\left[ Q_a, Q_b \right] = i f^c_{ab} Q_c + i \kappa_{ab} \mathbb{I}\)
    \item \(\kappa_{ab} = \omega_{ab} - \omega_{ba}\)
    \item \([Q_a,Q_b] + [Q_b,Q_a]=0 \implies \kappa_{ab}+ \kappa_{ba}=0\)
    \item \([Q_a,[Q_b,Q_c]] + [Q_b,[Q_c,Q_a]] + [Q_c,[Q_a,Q_b]]=0\) (Jacobi identity)
\end{itemize}

The structure constants respects those identity in the following form
\begin{itemize}
    \item \(f^a_{bc} + f^a_{cb} = 0\)
    \item \(f^d_{ab} f^e_{dc} + f^d_{bc} f^e_{da} + f^d_{ca} f^e_{db} = 0\) (Jacobi identity)
\end{itemize}
and along with the central charges they respect the following identity
\[
    \kappa_{ad} f^d_{bc} + \kappa_{bd} f^d_{ca} + \kappa_{cd} f^d_{ab} = 0,
\]
which is a consequence of the Jacobi identity for the generators \(Q_a\) of the representation \(U\) of the Lie group \(G\). A set of coefficients \(f^c_{ab}\) and \(\kappa_{ab}\) which respect the above identities is called a \textbf{Lie algebra} \(\mathfrak{g}\) \textbf{2-cocycle}. This naming is consistent since the central charges \(\kappa_{ab}\) are related to the phase factor \(\gamma\) of the projective representation \(U\) of the Lie group \(G\) by the relation \(\kappa_{ab} = \omega_{ab} - \omega_{ba}\).

...
\[
    U(g,g^{\prime}) = \gamma(g,g^{\prime}) U(g)U(g^{\prime}),
\]
\[
    U(g) \to \hat{U}(g) = \alpha(g) U(g),
\]
\[
    \gamma(g,g^{\prime}) = \alpha(g, g^{\prime}) \alpha(g)^{-1} \alpha(g^{\prime})^{-1},
\]
so that we can redefine the unitary operators \(U(g)\) of the projective representation \(U\) of the Lie group \(G\) by a suitable phase factor \(\alpha(g)\) to get a new projective representation \(\hat{U}\) of the Lie group \(G\) with a new phase factor \(\hat{\gamma}\) which is related to the old one by the relation
\[
    \hat{\gamma}(g,g^{\prime}) = \alpha(g, g^{\prime}) \alpha(g)^{-1} \alpha(g^{\prime})^{-1} \gamma(g,g^{\prime}),
\]
so that if we can find a suitable phase factor \(\alpha(g)\) such that \(\hat{\gamma}(g,g^{\prime}) = 1\) for all \(g, g^{\prime} \in G\), then we can redefine the unitary operators \(U(g)\) of the projective representation \(U\) of the Lie group \(G\) by a suitable phase factor \(\alpha(g)\) to get a new projective representation \(\hat{U}\) of the Lie group \(G\) with a trivial phase factor \(\hat{\gamma}\), so that the new projective representation \(\hat{U}\) of the Lie group \(G\) is actually a linear representation of the Lie group \(G\). However, in general, it is not always possible to find such a phase factor \(\alpha(g)\), so that there are some projective representations of the Lie group \(G\) which cannot be redefined to get a linear representation of the Lie group \(G\).

???

\[
    Q_a \to \hat{Q}_a = Q_a + k_a \mathbb{I},
\]
where \(k_a\) are the central shifts of the generators \(Q_a\) of the representation \(U\) of the Lie group \(G\). These shifts are related to the central charges \(\kappa_{ab}\) by the relation
\[
    \kappa_{ab} = f^c_{ab} k_c,
\]
which defines a \textbf{2-coboundary} Lie algebra. A 2-coboundary Lie algebra is special case of a 2-cocycle Lie algebra, where the central charges \(\kappa_{ab}\) can be expressed as a linear combination of the structure constants \(f^c_{ab}\) with coefficients given by the central shifts \(k_a\).
This is \textbf{Lie Algebra Cohomografy}

\subsection{Exponential Parametrization}

Starting from the exponential parametrization of the Lie group \(G\)
\[
    \tilde{g}(\underline{\theta}) = \lim_{n \to \infty} g(\frac{\theta}{n})^n = e^{\theta^a T_a},
\]
while for the inverse and the multiplication we have
\[
    \begin{dcases}
        \tilde{j}(\underline{\theta}) = - \theta^a, \\
        \tilde{m}(\underline{\theta}, \underline{\theta}') = \underline{\theta} + \underline{\theta}' + \frac{1}{2} f_{ab}^c \theta^a \theta'^b + \mathrm{O}(\theta^3),
    \end{dcases}
\]
where the last function is the Baker-Campbell-Hausdorff formula (BCH) for the multiplication of the Lie group \(G\) in the exponential parametrization.

We can now use this exponential parametrization of the Lie group \(G\) to get an exponential parametrization of the projective representation \(U\) of the Lie group \(G\):
\[
    \tilde{U}(g(\underline{\theta})) = U(\tilde{g}(\underline{\theta})) = U(e^{\theta^a T_a}) = 1- i \theta^a Q_a + \frac{1}{2} \theta^a \theta^b Q_{ab} + \mathrm{O}(\theta^3).
\]
For the phase function \(\omega\) we have
\[
    \tilde{\omega}(\underline{\theta}, \underline{\theta}') = \tilde{\omega}_{ab} \theta^a \theta'^b + \mathrm{O}(\theta^3).
\]
At the linear level there are no differences between \(g\) and \(\tilde{g}\), they only manifest from the second order in the parameters \(\theta^a\). Thus we can derive
\[
    Q_a = \tilde{Q}_a,
\]
and for the phase function \(\omega\) we have
\[
    \tilde{Q}_{ab} = \frac{1}{2}( (\omega_{ab} + \omega_{ba}) \mathbb{I} - Q_a Q_b - Q_b Q_a).
\]

Now we are interested in knowing if the following is true:
\[
    \tilde{U}(\theta) = \exp(i\theta^a Q_a),
\]
which is not the case: this is a projected representation, so this will be proven to be true only up to a phase factor. We can start from the exponential parametrization of the projective representation \(U\) of the Lie group \(G\)
\[
    \tilde{U}(\underline{\theta}) = U(\tilde{g}(\underline{\theta})) = \lim_{n \to \infty} U(g(\frac{\theta}{n}))^n = \lim_{n \to \infty} U((\frac{1}{n^2})^n) \neq \lim_{n \to \infty} U((\frac{1}{n^2}))^n,
\]
where we highlighted the difference between this projected and a true representation. Computing further we get
\[
    \tilde{U}(\underline{\theta}) = \exp(i \sum_{k}^{n} \omega(\frac{\theta}{n}, \frac{k\theta}{n})) U(\frac{\underline{\theta}}{n})^n,
\]
where the last term, as we take the limit \(n \to \infty\), becomes
\[
    \lim_{n \to \infty} U(\frac{\underline{\theta}}{n})^n = (1 - i \frac{\theta^a}{n} Q_a)^n = e^{-i \theta^a Q_a},
\]
so that the exponential parametrization of the projective representation \(U\) of the Lie group \(G\) is given by
\[
    \tilde{U}(\underline{\theta}) = e^{i \alpha(\underline{\theta})} e^{-i \theta^a Q_a}.
\]

If we now want to get .....

\[
    e^{i \alpha(\underline{\theta}) - i \theta^a k_a} \tilde{U}(\underline{\theta}) = e^{-i \theta^a \hat{Q}_a},
\]
where \(\hat{Q}_a = Q_a + k_a \mathbb{I}\) are the shifted generators of the representation \(U\) of the Lie group \(G\): these operates in the same way since they differ only by a phase factor, but they have different central charges.

Now we found a true representation
\begin{itemize}
    \item \(\hat{\tilde{U}}(\tilde{\underline{j}}(\underline{\theta})) = \hat{\tilde{U}}(\underline{\theta})^{-1}\)
    \item \(\hat{\tilde{U}}(\tilde{\underline{m}}(\underline{\theta}, \bs{\theta'})) = \hat{\tilde{U}}(\underline{\theta}) \hat{\tilde{U}}(\bs{\theta'})\)
\end{itemize}

The term \(\exp(-i \theta^a k_a)\) is an important \textbf{phase factor}, since it is the one which allows us to get a true representation from a projective representation by shifting the generators \(Q_a\) of the representation \(U\) of the Lie group \(G\) by a suitable central shift \(k_a\). This is the reason why we introduced the concept of 2-coboundary Lie algebra, since it allows us to understand when a projective representation of a Lie group can be redefined to get a true representation of the Lie group.

This is a \textbf{multi-valued function} on the underlying group manifold \(G\), since the phase factor usually is not an integer which comes back to the original value after a loop around the group manifold \(G\) (it is not guaranteed that the passage is single-valued): the exponential parametrization of the projective representation \(U\) of the Lie group \(G\) is a multi-valued function on the underlying group manifold \(G\). Usually then we introduce a branch cut in the complex plane to make it single-valued, in order to treat multi valued function as single valued function on the more complex Riemann surface (such as for the square root and the logarithm).

\[
    \ln z = \ln \vert z \vert  + i \arg z + 2 \pi i n, \quad n \in \mathbb{Z}.
\]

In our case the exponential parametrization of the projective representation \(U\) of the Lie group \(G\) is a multi-valued function on the underlying group manifold \(G\), but can be made single-valued on a infinite Riemann surface which is a covering space of the underlying group manifold \(G\): a covering group \(\tilde{G}\). It is the same idea of the Riemann surface, but with more subtleties and complications, since the group structure has to be preserved.

There can be made analogies with the various covering groups we can find in physics, such as the universal covering group of the Lorentz group, which is the Spin group, and the universal covering group of the rotation group \(SO(3)\), which is the special unitary group \(SU(2)\). The projective representations of the Lorentz group and the rotation group \(SO(3)\) can be redefined to get true representations of their respective covering groups, which are the Spin group and the special unitary group \(SU(2)\). This is a consequence of the fact that the central charges of the projective representations of the Lorentz group and the rotation group \(SO(3)\) can be expressed as linear combinations of the structure constants of their respective Lie algebras with coefficients given by suitable central shifts
\[
    \mathrm{SL}(2, \mathbb{C}) = \widetilde{\mathrm{SO}^+(1,3)}, \quad \mathrm{SU}(2) = \widetilde{\mathrm{SO}(3)}.
\]

\begin{example}[The translation group \(\mathrm{T}(3)\)]
    The translation group \(\mathrm{T}(3)\) implements the translation symmetry of the three-dimensional space:
    \[
        \mathbf{a} \in \mathbb{E}_3, \quad T(\mathbf{x}) = \mathbf{x} + \mathbf{a},
    \]
    where \(T \in \mathrm{T}(3)\) and \(\mathbb{E}_3\) is the three-dimensional Euclidean space.

    The proper translation group \(\mathrm{T}(3)\) can be represented by the unitary operators \(U(\mathbf{a})\) for a translation \(\mathbf{a} \in \mathbb{E}_3\) which act on the Hilbert space of states of a quantum system as
    \[
        (U(\mathbf{a})\Psi)(\mathbf{x}) = \Psi(\mathbf{x} + \mathbf{a}),
    \]
    where \(\Psi\) is a state of the quantum system and we want to prove that \(\{U(\mathbf{a})\}_{\mathbf{a} \in \mathbb{E}_3}\) yields a unitary representation of the proper translation group \(\mathrm{T}(3)\).
    \begin{proof}
        We can start by computing the composition of two translations \(\mathbf{a}\) and \(\mathbf{b}\):
        [...]

        There are no multiplicative phases in the right hand side. The operator family \(U(\mathbf{a})\) constitutes in this way a unitary representation of the proper translation group \(\mathrm{T}(3)\).
    \end{proof}
    Let us now discuss the infinitesimal generators of the translation group \(\mathrm{T}(3)\). The infinitesimal generators \(\hbar^{-1}\mathbf{p}\) of the unitary representation \(U\) of the proper translation group \(\mathrm{T}(3)\) are defined through the expansion of the unitary operators \(U(\mathbf{a})\) for small translations \(\delta\mathbf{a}\):
    \[
        U(\delta\mathbf{a}) = 1 + \frac{i}{\hbar} \mathbf{p} \cdot \delta\mathbf{a} + \mathrm{O}(\delta a^2),
    \]
    with \(\mathbf{p} = (p_1, p_2, p_3)\) being the momentum operator of the quantum system
    \[
        \mathbf{p} \Psi(\mathbf{x}) = -i \hbar \nabla \Psi(\mathbf{x}).
    \]
    Thus we can study the expansion of the unitary operators \(U(\mathbf{a})\) acting on a state \(\Psi\) for small translations \(\delta\mathbf{a}\):
    \[
        \begin{aligned}
            (U(\delta\mathbf{a})\Psi)(\mathbf{x}) & = \Psi(\mathbf{x} + \delta\mathbf{a})                                                         \\
                                                  & = \Psi(\mathbf{x}) + \delta \mathbf{a} \cdot \nabla \Psi(\mathbf{x}) + \mathrm{O}(\delta a^2) \\
                                                  & = \dots
        \end{aligned}
    \]

    \(\mathbf{p}\) is therefore the familiar quantum linear momentum operator as anticipated by the notation used. The result obtained is in line with the standard interpretation of momentum as the infinitesimal generator of configuration space translations.

    As well–known from quantum theory, its components
    \[
        p_k = \mathbf{p} \cdot \mathbf{e}_k, \quad k = 1, 2, 3,
    \]
    satisfy the commutation relations
    \[
        [p_k, p_l] = 0, \quad k, l = 1, 2, 3,
    \]
    free of central charges, as expected from the fact that the translation group \(\mathrm{T}(3)\) is an abelian group, so that its Lie algebra is also abelian, with vanishing structure constants \(f^c_{ab} = 0\).

    The translation group \(\mathrm{T}(3)\) is a special case of a more general class of groups called the Heisenberg group, which is a non-abelian group with non-vanishing structure constants and non-vanishing central charges. The Heisenberg group is the symmetry group of the quantum phase space, and its projective representations are related to the canonical commutation relations of quantum mechanics.

    Thus the operators \(U(\mathbf{a})\) can be written in terms of the momentum operator \(\mathbf{p}\), the generators, as
    \[
        U(\mathbf{a}) = e^{\tfrac{i}{\hbar} \mathbf{a} \cdot \mathbf{p}}.
    \]
    \begin{proof}
        limit for n....
    \end{proof}


    .


\end{example}

\begin{example}[The proper rotation group \(\mathrm{SO(3)} / \mathrm{SU}(2)\)]

    The proper rotation group \(SO(3)\) implements the rotation symmetry of the three-dimensional space:
    \[
        R \in SO(3), \quad R: \mathbf{x} \to R\mathbf{x},
    \]
    where \(R\) is a proper rotation and \(\mathbf{x}\) is a vector in the three-dimensional space.

    The proper rotation group \(SO(3)\) can be represented by the unitary operators \(U(R)\) for a proper rotation \(R \in SO(3)\) which act on the Hilbert space of states of a quantum system as
    \[
        (U(R)\Psi)(\mathbf{x}) = D(\mathbf{R})\Psi(R^{-1}\mathbf{x}),
    \]
    where \(\Psi\) is a state of the quantum system, \(\mathbf{x}\) is a point in the configuration space, and \(D(\mathbf{R})\) is a unitary operator which acts on the internal degrees of freedom of the quantum system, such as spin.

    Indeed we know that the spin of a quantum system is related to the representation of the rotation group \(SO(3)\) on the internal degrees of freedom of the quantum system; this does not come as a surprise, since the spin of a quantum system is related to the way the quantum system transforms under rotations, and the rotation group \(SO(3)\) is the symmetry group of rotations in three-dimensional space.

    We parametrized the representation in order to act on \(\mathbf{x}\) as the inverse of the rotation,\footnote{Which is always well defined from the properties of groups.} since this assumption is consistent with the composition of rotations, as we can see by computing the composition of two rotations \(R\) and \(R'\):
    \[
        U(\mathbf{R},\mathbf{R}') = U(\mathbf{R}) U(\mathbf{R}'),
    \]
    otherwise the the order would be inverted.

    \(D(\mathbf{R})\) are \((2s+1)\times(2s+1)\) unitary matrices which form a representation of the rotation group \(SO(3)\) on the internal degrees of freedom of the quantum system, where \(s\) is the spin of the quantum system
    \[
        D(\mathbf{R})^{\dagger}D(\mathbf{R}) = D(\mathbf{R}) D(\mathbf{R})^{\dagger} = \mathbb{I},
    \]
    and they satisfy the composition law\footnote{Here the right order is implemented by the previously discussed choice of parametrization.}
    \[
        D(\mathbf{R} \mathbf{R}') = D(\mathbf{R}) D(\mathbf{R}').
    \]
    \begin{proof}
        Let us show first that the operator \(U(R)\) is unitary. Define the operator
        \[
            U^*(\mathbf{R})\Psi(\mathbf{x}) = D(\mathbf{R})^\dagger \Psi(R\mathbf{x}),
        \]
        so that we can compute the composition of \(U^*(\mathbf{R})\) and \(U(\mathbf{R})\) which acts as the identity operator on the Hilbert space of states of the quantum system:
        \[
            (U^*(\mathbf{R}) U(\mathbf{R})\Psi)(\mathbf{x}) = D(\mathbf{R})^\dagger D(\mathbf{R}) \Psi(\mathbf{R}^{-1}\mathbf{R}\mathbf{x}) = \Psi(\mathbf{x}).
        \]
        The next thing to show is that
        \[
            \Vert U(\mathbf{R})\Psi \Vert^2 = \Vert \Psi \Vert^2,
        \]
        which is the definition of unitarity. We can compute
        \[
            \begin{aligned}
                \Vert U(\mathbf{R})\Psi \Vert^2 & = \int d^3x \, |U(\mathbf{R})\Psi(\mathbf{x})|^2 = \int d^3x \, |D(\mathbf{R})\Psi(R^{-1}\mathbf{x})|^2 \\
                                                & = \int d^3x \, |\Psi(R^{-1}\mathbf{x})|^2 = \int d^3x' \, |\Psi(\mathbf{x}')|^2 = \Vert \Psi \Vert^2,
            \end{aligned}
        \]
        where we have used the unitarity of \(D(\mathbf{R})\) and \(J = |\det(\mathbf{R})^{-1}| = 1\). The last step is to show that the composition law is satisfied, which can be done by computing
        \[
            \begin{aligned}
                (U(\mathbf{R} \mathbf{R}')\Psi)(\mathbf{x}) & = D(\mathbf{R} \mathbf{R}')\Psi((\mathbf{R} \mathbf{R}')^{-1}\mathbf{x}) = D(\mathbf{R}) D(\mathbf{R}')\Psi(\mathbf{R}'^{-1} \mathbf{R}^{-1}\mathbf{x}) \\
                                                            & = D(\mathbf{R}) (U(\mathbf{R}')\Psi)(\mathbf{R}^{-1}\mathbf{x}) = (U(\mathbf{R}) U(\mathbf{R}')\Psi)(\mathbf{x}),
            \end{aligned}
        \]
        proving that \(U(\mathbf{R} \mathbf{R}') = U(\mathbf{R}) U(\mathbf{R}')\).
    \end{proof}

    We now aim to study the \textbf{exponential parametrization} of the representation \(U\) of the proper rotation group \(SO(3)\). The generators \(J_k\) of the representation \(U\) of the proper rotation group \(SO(3)\) are defined through the expansion of the unitary operators \(U(R)\) for small rotations \(\delta R\).

    We can then parametrize the rotation \(R\) by a vector \(\bs{\varphi} = (\varphi_1, \varphi_2, \varphi_3)\) which represents the axis of rotation and the angle of rotation, so that the exponential parametrization of the the proper rotation group \(SO(3)\) is given by
    \[
        R(\bs{\varphi}) = e^{- * \varphi}, \quad R(\bs{\varphi})^{-1} = e^{* \varphi},
    \]
    where
    \[
        *\varphi = \begin{pmatrix}
            0          & \varphi_3  & -\varphi_2 \\
            -\varphi_3 & 0          & \varphi_1  \\
            \varphi_2  & -\varphi_1 & 0
        \end{pmatrix},
    \]
    is the antisymmetric matrix associated to the vector \(\bs{\varphi} = (\varphi_1, \varphi_2, \varphi_3)\) which parametrizes the rotation \(R\).

    Now we can reflect this exponential parametrization on the representation \(U\) of the proper rotation group \(SO(3)\) to get an exponential parametrization of the representation:
    \[
        (U(\bs{\varphi})\Psi)(\mathbf{x}) = D(\bs{\varphi})\Psi(e^{* \varphi}\mathbf{x}),
    \]
    where \(D(\bs{\varphi})\) is the exponential parametrization of the representation of the rotation group \(SO(3)\) on the internal degrees of freedom of the quantum system. We have to look at the expansion of the unitary operators \(U(\bs{\varphi})\) for small rotations \(\delta \bs{\varphi}\) to get the generators \(\underline{j}\) of the representation \(U\) of the proper rotation group \(SO(3)\):
    \[
        U(\delta \bs{\varphi}) = 1 - i \underline{j} \cdot \delta \bs{\varphi} + \mathrm{O}(\delta \varphi^2),
    \]
    where \(\underline{j} = (J_1, J_2, J_3)\) are the generators of the representation \(U\) of the proper rotation group \(SO(3)\), which are called the angular momentum operators of the quantum system. Then the expansion of the internal degrees transformation \(D(\bs{\varphi})\) for small rotations \(\delta \bs{\varphi}\) gives us
    \[
        D(\delta \bs{\varphi}) = 1 - i \bs{\Sigma} \cdot \delta \bs{\varphi} + \mathrm{O}(\delta \varphi^2),
    \]
    where \(\bs{\Sigma} = (\Sigma_1, \Sigma_2, \Sigma_3)\) are the generators of the representation of the rotation group \(SO(3)\) on the internal degrees of freedom of the quantum system, which are called the spin operators of the quantum system. Finally we have to compute the expansion of the spatial transformation \(e^{* \varphi}\) for small rotations \(\delta \bs{\varphi}\) to get
    \[
        e^{* \delta \varphi} = 1 + *\delta \varphi + \mathrm{O}(\delta \varphi^2),
    \]
    but since the rotation can be expressed as \(- \delta \bs{\varphi} \times \mathbf{x}\), we can write
    \[
        e^{* \delta \varphi} \mathbf{x} = \mathbf{x} - \delta \bs{\varphi} \times \mathbf{x} + \mathrm{O}(\delta \varphi^2).
    \]

    Now we can put everything together to get the expansion of the unitary operators \(U(\bs{\varphi})\) for small rotations \(\delta \bs{\varphi}\):
    \[
        \begin{aligned}
            (U(\delta \bs{\varphi})\Psi)(\mathbf{x}) & = D(\delta \bs{\varphi})\Psi(e^{* \delta \varphi}\mathbf{x})                                                                                                          \\
                                                     & = (1 - i \bs{\Sigma} \cdot \delta \bs{\varphi} + \mathrm{O}(\delta \varphi^2))\Psi(\mathbf{x} - \delta \bs{\varphi} \times \mathbf{x} + \mathrm{O}(\delta \varphi^2)) \\
                                                     & = (1 - i \bs{\Sigma} \cdot \delta \bs{\varphi} + i (\delta \bs{\varphi} \times \mathbf{x}) \cdot (-i\nabla) + \mathrm{O}(\delta \varphi^2))\Psi(\mathbf{x})           \\
                                                     & \dots                                                                                                                                                                 \\
                                                     & = \Psi(\mathbf{x}) - \delta \bs{\varphi} \cdot (\mathbf{x} \times \nabla + i \bs{\Sigma})\Psi(\mathbf{x}) + \mathrm{O}(\delta \varphi^2),
        \end{aligned}
    \]
    [...]

    Thus we found
    \[
        \underline{j} = \hbar \bs{\Sigma} - i\hbar \mathbf{x} \times \nabla,
    \]
    where renaming \(\hbar \bs{\Sigma} = \mathbf{S}\) and \(-i \hbar \mathbf{x} \times \nabla = \mathbf{L}\) we get the familiar decomposition of the angular momentum operator \(\underline{j}\) into the spin operator \(\mathbf{S}\) and the orbital angular momentum operator \(\mathbf{L}\):
    \[
        \underline{j} = \mathbf{S} + \mathbf{L}.
    \]

    We know that the infinitesimal generators of a continuous symmetry group satisfy the same commutation relations as the Lie algebra of the symmetry group, so that the generators \(J_k\) of the representation \(U\) of the proper rotation group \(SO(3)\) satisfy the commutation relations
    \[
        [J_k, J_l] = i \hbar \epsilon_{klm} J_m, \quad k, l, m = 1, 2, 3,
    \]
    while the generators \(\Sigma_k\) of the representation of the rotation group \(SO(3)\) on the internal degrees of freedom of the quantum system satisfy the same commutation relations
    \[
        [\Sigma_k, \Sigma_l] = i \epsilon_{klm} \Sigma_m, \quad k, l, m = 1, 2, 3,
    \]
    where \(\epsilon_{klm}\) is the Levi-Civita symbol. The commutation relations of the generators \(J_k\) and \(\Sigma_k\) of the representation, are in line with the structure constants of the Lie algebra \(\mathfrak{so}(3)\) of the proper rotation group \(SO(3)\), which are given by
    \[
        f^m_{kl} = \epsilon_{klm}.
    \]

    Since we are using the exponential parametrization \(\mathbf{R}(\bs{\varphi})\) of \(\mathrm{SO}(3)\), we can get the exponential parametrization of the representation \(U(\bs{\varphi})\) in terms of the infinitesimal generators \(J_k\):
    \[
        U(\bs{\varphi}) = e^{-\tfrac{i}{\hbar} \bs{\varphi} \cdot \underline{j}}.
    \]
    \begin{proof}
        limit for n....
    \end{proof}

\end{example}


It should be made an extensive digression on the universal covering group, which is the group which allows us to get a single-valued representation from a multi-valued representation by lifting the representation to the covering group. The universal covering group of a Lie group \(G\) is a simply connected Lie group \(\tilde{G}\) which covers \(G\) in such a way that the projection map from \(\tilde{G}\) to \(G\) is a local diffeomorphism. The universal covering group \(\tilde{G}\) has the same Lie algebra as \(G\), but it can have different global properties, such as different topology and different fundamental group. The universal covering group \(\tilde{G}\) allows us to get a single-valued representation from a multi-valued representation of \(G\) by lifting the representation to \(\tilde{G}\).

However we do not have the time to make this digression, so we will just mention that the universal covering group of the proper rotation group \(SO(3)\) is the special unitary group \(SU(2)\), which is a simply connected Lie group that covers \(SO(3)\) in such a way that the projection map from \(SU(2)\) to \(SO(3)\) is a local diffeomorphism. The universal covering group \(SU(2)\) has the same Lie algebra as \(SO(3)\), but it has different global properties, such as different topology and different fundamental group. The universal covering group \(SU(2)\) allows us to get a single-valued representation from a multi-valued representation of \(SO(3)\) by lifting the representation to \(SU(2)\).